%==============================================================================
% Capítulo 9: Conclusiones
%==============================================================================
\section{Conclusiones}\label{sec:conclusiones}

El trabajo desarrollado ha cumplido con el objetivo general planteado en la \autoref{subsec:objetivo-general}: construir un sistema automatizado de mesa de ayuda capaz de recibir, procesar, clasificar y registrar incidentes de manera automática, reduciendo la intervención humana en la etapa inicial. La integración de N8N como motor de orquestación, un módulo Python con FastAPI como núcleo de procesamiento, PostgreSQL como capa de persistencia, Twilio como canal telefónico y el modelo Gemini~3.6 Flash como motor de inferencia lingüística ha demostrado ser una combinación arquitectónica viable, eficiente y económicamente accesible para organizaciones medianas. Ese cumplimiento del objetivo general no se extiende al segundo objetivo específico: la exactitud estricta del 73,5~\% y el F1 macro quedaron por debajo de las cotas mínimas de 85~\% y 0,85 fijadas en la \autoref{subsec:objetivos-especificos}, de modo que esa meta se declara no alcanzada. En cambio, la reducción del tiempo medio de registro del 71,2~\% sí se sostiene como resultado central de la validación experimental del sistema, con significación estadística; el detalle de las métricas y del contraste se reporta en la \autoref{sec:resultados}. En consecuencia, la hipótesis principal se declara parcialmente confirmada ---en la reducción del tiempo de registro--- y parcialmente refutada ---en el umbral de clasificación---.

\subsection{Cumplimiento de los objetivos del trabajo}\label{subsec:cumplimiento-objetivos}

El diseño arquitectónico del sistema materializó los principios de separación de responsabilidades y cohesión interna que la disciplina de ingeniería de software prescribe para sistemas distribuidos. La arquitectura de cinco capas ---canales de entrada, orquestación, procesamiento, inferencia y persistencia--- distribuyó las responsabilidades entre componentes independientes comunicados mediante interfaces REST. La elección de tecnologías maduras y de despliegue autoalojado (Docker, PostgreSQL, N8N) no fue meramente una decisión instrumental sino una respuesta directa a las restricciones del contexto: organizaciones medianas argentinas que operan con presupuestos acotados, enfrentan limitaciones cambiarias para la adquisición de software comercial y están sujetas a un marco normativo que exige control sobre los datos personales.

Sobre esa arquitectura se construyó el clasificador híbrido, cuyo desempeño sobre el corpus real quedó por debajo de las cotas mínimas fijadas en el segundo objetivo específico, por lo que ese objetivo no se alcanza; la brecha y sus valores se analizan en la \autoref{sec:discusion}. La estrategia compositiva ---un filtro determinístico de primera etapa y un modelo de lenguaje grande reservado para los casos que el filtro no resuelve con confianza suficiente--- resultó, sin embargo, el rasgo más relevante del sistema: la mayor parte del volumen se resuelve sin inferencia externa y el modelo se emplea como árbitro de último recurso antes que como clasificador universal, con el equilibrio entre costo, latencia y exactitud que se detalla en la \autoref{subsec:subsistema-clasificacion}.

La integración de los tres canales paralelos dentro de un único flujo de orquestación confirmó que la diversidad de vías de entrada no introduce dispersión operativa cuando la normalización se realiza tempranamente en el flujo. El canal telefónico, que a priori representaba el mayor desafío técnico por la inclusión de transcripción automática del habla, alcanzó un tiempo de registro medio de 20,1 segundos, que incluye la duración de la llamada grabada además de la latencia de procesamiento ---un valor que lo posiciona como una vía prácticamente útil y no como una demostración de laboratorio. La decisión de normalizar la entrada en el propio N8N, antes de invocar al servicio de clasificación, resultó acertada: aisló la complejidad de cada protocolo de recepción del núcleo de procesamiento, permitiendo que el servicio Python operara siempre sobre la misma estructura de datos independientemente del canal de origen.

La evaluación comparativa entre el flujo manual y el automatizado se presenta en la \autoref{subsec:comparacion-tiempos} y se interpreta, con sus alcances y límites, en la \autoref{subsec:interpretacion-resultados}. En síntesis, el sistema automatizado resultó más rápido que el operador humano en 188 de los 200 pares medidos, con la reducción del tiempo de registro como resultado central. La intervención humana por caso en el flujo automatizado no se midió, de modo que el aporte operativo asociado al gate de revisión se declara como no medido: el sistema libera al personal técnico de la carga rutinaria de registro, redirigiendo ese tiempo hacia la resolución efectiva de los casos. El reporte de los estadísticos de la prueba y su discusión no se repiten aquí.

La documentación generada ---arquitectura, procedimientos de despliegue, especificación OpenAPI, flujo N8N exportable y análisis tipificado de errores--- constituye el aporte documental del quinto objetivo específico. La inclusión de una tipología de errores con tres patrones identificados (descripciones cortas, mezcla de categorías y jerga local) proporciona un punto de partida concreto para cualquier organización que desee replicar o adaptar el sistema, permitiéndole anticipar las categorías de fallo más probables antes de iniciar su propia implementación.

\subsection{Aportes y generalización}\label{subsec:aportes-generalizacion}

Más alla del cumplimiento de los objetivos específicos, los resultados de este trabajo permiten formular una conclusión acotada a este corpus y al contraste efectivamente realizados: en este corpus de 200 incidentes de mesa de ayuda, en español rioplatense y con cinco categorías canónicas, un modelo de lenguaje grande invocado mediante instrucciones contextualizadas ---sin ajuste fino sobre datos del dominio--- resultó una alternativa viable para la clasificación de tickets, siempre que se combine con un filtro determinístico de primer nivel que absorba los casos triviales y reduzca el costo de inferencia. Esta lectura se limita a este dominio: no se realizó una comparación directa contra clasificadores supervisados entrenados sobre el mismo corpus, de modo que no se sostiene aquí una superioridad general de los modelos de lenguaje frente a esos clasificadores (véase la \autoref{subsec:limitaciones-estudio}).

Esta combinación arquitectónica ---orquestación visual, filtrado determinístico y LLM como árbitro semántico--- constituye un patrón de diseño reproducible para el dominio de gestión de servicios en organizaciones medianas de América Latina, donde la restricción presupuestaria y la soberanía de datos son determinantes para la adopción tecnológica. La arquitectura propuesta no depende de infraestructura especializada de GPU ni de licencias de software propietario; su costo operativo está dominado por el consumo de API de inferencia, cuyo orden de magnitud y su relación con el ahorro de tiempo se cuantifican en la \autoref{sec:discusion}.

El trabajo aporta tres contribuciones que se sostienen más alla del caso particular. La primera es una arquitectura reproducible documentada en todos sus niveles, desde el modelo de datos hasta el contrato REST y el flujo de orquestación, que cualquier equipo con conocimientos básicos de Docker, Python y N8N puede replicar. La segunda es evidencia empírica, recolectada bajo un diseño experimental riguroso con doble etiquetado independiente y prueba no paramétrica, sobre el desempeño de LLM en un dominio y una variante lingüística escasamente documentados en la literatura. La tercera es un marco de implementación que contempla desde el inicio las restricciones normativas argentinas, ofreciendo un camino de adopción que no requiere sortear barreras legales ni asumir riesgos de cumplimiento.

\newpage
